\documentclass{article}

%%%%%%%%%%%%%%% LIBRERIAS %%%%%%%%%%%%%%%%%%%%%
\usepackage{amsmath}
\usepackage{titlesec}
\usepackage{titletoc}
\usepackage{graphicx}
\usepackage[spanish,es-tabla]{babel} % 'es-tabla' cambia Cuadro→Tabla
\usepackage{hyperref}                % cargar después de babel
\usepackage{float}
\usepackage{circuitikz}
%%%%%%%%%%%%%%%%%%% VARIABLES %%%%%%%%%%%%%%%%%%%%
\newcommand{\Facultad}{Instituto Tecnológico \\de\\ Buenos Aires} %constantes
\newcommand{\TPn}{Trabajo Práctico N° 3}
\newcommand{\TPtema}{Respuesta Transitoria}
\renewcommand{\thesection}{\arabic{section}}          % 2
\renewcommand{\thesubsection}{\quad \alph{subsection}}   % a
\renewcommand{\thesubsubsection}{\quad \thesubsection.~\roman{subsubsection}} % a. i
\graphicspath{{imagenes/}} %para que acceda a las fotos en la carpeta directamente


%%%%%%%%%%%%%%%%%% FORMATO TÍTULO Y NUMERACIÓN %%%%%%%%%%%%%%%%%%%
% Numeración de secciones
\renewcommand{\thesection}{\arabic{section}.}          
\renewcommand{\thesubsection}{\thesection\arabic{subsection}}       
\renewcommand{\thesubsubsection}{$\alph{subsubsection})$}

% Numerar hasta subsubsecciones
\setcounter{secnumdepth}{3}

% Formato de títulos
\titleformat{\section}{\Huge\bfseries}{\thesection}{1em}{}
\titleformat{\subsection}{\LARGE\bfseries}{\thesubsection}{0.5em}{}
\titleformat{\subsubsection}{\large\bfseries}{\thesubsubsection}{0.5em}{}



%%%%%%%%%%%%%%%%%%% ARCHIVO %%%%%%%%%%%%%%%%%%%%%%%%
\begin{document}

%%%CARATULA%%%
\begin{titlepage} %creo portada

        \begin{flushleft}
            \centering
            \includegraphics[width=0.3\textwidth]{Logo_ITBA.png}
        \end{flushleft}

        \centering
            
        {\scshape\LARGE \Facultad \par} %\par sirve para indicar un final de parrafo
        \vspace{1cm}                    %esto hace un espacio entre lineas de 1cm


        {\huge\bfseries \TPn \par}
        \vspace{1.5cm}
        {\Large Teoría de Circuitos I\\ 25.10 \par}
        \vfill                      %sirve para rellenar el espacio y quede simétrico. Si se añaden otros, se dividen el espacio de forma equitativa
        {\Large \bfseries Grupo N° 5 \par}
        \vspace{1cm}
        {\large Juan Bautista Correa Uranga \hfill Legajo: 65016 \par} %\hfill sirve para hacerlo simétrico
        {\large Juan Ignacio Caorsi \hfill Legajo: 65532  \par}
        {\large Rita Moschini \hfill Legajo: 67026 \par} 
        \vfill
        {\large \today\par}
        \vfil

    \end{titlepage}

 %%%RESUMEN%%%
{\centering \LARGE \bfseries Resumen \par}

\newpage

%%%INDICE%%%
\tableofcontents %esto sirve para crear el índice

\newpage

%%%Introduccion%%%
\section{Introducción}
    \subsection{Instrumental}
    \subsection{Marco teórico}

%%%Desarrollo%%%
\section{Desarrollo}
    \subsection{Procedimiento}
    \subsection{Mediciones}
 
        \begin{itemize}
            \item $ R_f = 215 \Omega $
            \item $ R_{V_{max}} = 9980 \Omega $
            \item $ R_{V_{min}} = 2 \Omega $
            \item $ R_L = 0,8 \Omega $
            \item $ L \approx 1 mH $
        \end{itemize}
        
        \par
        \subsubsection*{Capacitor de C = 470 pF}
            \begin{itemize}
                \item Valor de la resistencia variable tal que el amortiguamiento era crítico: $ R_{critico} = 1,9 k\Omega $
                \item Tiempo $\tau$ en que la salida llegaba a 3,175 V cuando $ R_V $ tomaba su valor máximo: $ \tau = 5,75 \mu s$
                \item Salida cuando $t=5\tau$: V=3,175 V
            \end{itemize}

        \par
        \subsubsection*{Capacitor de C = 47 pF}
            \begin{itemize}
                \item Valor de la resistencia variable tal que el amortiguamiento era crítico: $ R_{critico} = 3,47 k\Omega $
                \item Tiempo $\tau$ en que la salida llegaba a 3,175 V cuando $ R_V $ tomaba su valor máximo: $ \tau = 32,25 \mu s$
                \item Tiempo en que la salida llegaba a 5,24 V ($ 5V \pm 0,05V $) cuando $ R_V $ tomaba su valor mínimo:  $ t = 14,3 \mu s$ %a frecuencia de oscilación y el sobrepico máximo
                \item Tiempo en que la salida llegaba a 4,888 V ($ 5V \pm 0,05V $) con ambas resistencias cortocircuitadas: $ t = 22,25 \mu s$
            \end{itemize}

    \subsection{Cálculos}

        Cálculo del valor de la resistencia variable tal que el amortiguamiento es crítico: \par
        $ \alpha_{serie} = \omega_0 $
        $ \Rightarrow \frac{R}{2L} = \frac{2L}{\sqrt{LC}} $
        $ \Rightarrow R = \frac{2L}{\sqrt{LC}}  \wedge R_f + R_V + R_L $
        $ \Rightarrow R_V = R - R_f  -R_L $

        Cálculo del valor de $ \tau $ para $ R_V = R_{V_{max}} = 9980 \Omega $
        $ \tau = \frac{1}{\alpha} = \frac{2L}{R_f + R_V + R_L} $

            \paragraph*{Capacitor de C = 470 pF} 
                \par Valor de la resistencia variable tal que el amortiguamiento fuera crítico: $ R_{critico} = 2,702 k\Omega $ %CUALQUIER COSA DESPUES PONEMOS MENOS DECIMALES!!! prefiero poner de mas para hacer menos veces la cuenta

            \paragraph*{Capacitor de C = 47 pF} 
                Valor de $ \tau $ para $ R_V = R_{V_{max} }= 9980 \Omega $: $ \tau = \frac{2 \cdot 10^-3}{215 + 9989 + 0,8} $ \par
                $ \tau = 0,91616 \mu s $
                \par Valor de la resistencia variable tal que el amortiguamiento fuera crítico: $ R_{critico} = 9,010 k\Omega $
        
        \subsubsection*{Resultados}
        \subsubsection*{Capacitor de C = 47 pF}


    \subsection{Análisis}
\section{Conclusiones}

\end{document}